\begin{definition}[Order Dense Subset]
\label{def:order-dense-subset}
Let $(S,\leq)$ be a linearly ordered set and let $D\subseteq S$.  The set
$D$ is \emph{order dense in} $S$ if between every two elements $x<y$ of
$S$ there lies an element of $D$.
\end{definition}
\end{definitionbox}

\begin{remark*}[Standard quantified statement]
\[
\begin{aligned}
&\operatorname{OrderDenseSubset}(D,P)
\Longleftrightarrow
\forall x,y\in S\;\bigl(x<y\Longrightarrow \exists d\in D\;(x<d<y)\bigr).
\end{aligned}
\]
\end{remark*}

\begin{remark*}[Predicate reading]
\[
S\in\mathbf{Set},\qquad D\subseteq S,\qquad P=\mathsf{TotallyOrderedSet}(S,\leq),\qquad
\operatorname{OrderDenseSubset}(D,P)\coloneqq
\forall x,y\in S\;\bigl(x<y\Longrightarrow \exists d\in D\;(x<d<y)\bigr).
\]
\end{remark*}

\begin{remark*}[Negated quantified statement]
\[
\begin{aligned}
&\text{Let }(S,\leq)\text{ be a linearly ordered set and }D\subseteq S.\\
&D\text{ is not order dense in }S
\Longleftrightarrow
\exists x,y\in S\;\bigl(x<y\land \forall d\in D\;(d\leq x\lor y\leq d)\bigr).
\end{aligned}
\]
\end{remark*}

\begin{remark*}[Negation predicate reading]
\[
S\in\mathbf{Set},\qquad D\subseteq S,\qquad P=\mathsf{TotallyOrderedSet}(S,\leq),\qquad
\neg\operatorname{OrderDenseSubset}(D,P)\Longleftrightarrow
\exists x,y\in S\;\bigl(x<y\land \forall d\in D\;(d\leq x\lor y\leq d)\bigr).
\]
\end{remark*}

\begin{remark*}[Failure modes]
\begin{description}
\item[Exposition.]
Order density fails when there is a nonempty open interval of $S$ that
contains no point of $D$.


  \item[\operatorname{OrderDenseSubset}.]
  \[
\begin{aligned}
&\text{Let }(S,\leq)\text{ be a linearly ordered set and }D\subseteq S.\\
&D\text{ is not order dense in }S
\Longleftrightarrow
\exists x,y\in S\;\bigl(x<y\land \forall d\in D\;(d\leq x\lor y\leq d)\bigr).
\end{aligned}
\]
\[
S\in\mathbf{Set},\qquad D\subseteq S,\qquad P=\mathsf{TotallyOrderedSet}(S,\leq),\qquad
\neg\operatorname{OrderDenseSubset}(D,P)\Longleftrightarrow
\exists x,y\in S\;\bigl(x<y\land \forall d\in D\;(d\leq x\lor y\leq d)\bigr).
\]
\end{description}
\end{remark*}

\begin{remark*}[Failure modes]
\begin{description}
\item[Exposition.]
\[
\neg\operatorname{OrderDenseSubset}(D,P)
=
\underbrace{\exists x,y\in S\;(x<y)}_{\text{nonempty interval}}
\;\land\;
\underbrace{\forall d\in D\;(d\leq x\lor y\leq d)}_{\text{no point of }D\text{ lies between }x\text{ and }y}.
\]


  \item[\operatorname{OrderDenseSubset}.]
  \[
\begin{aligned}
&\text{Let }(S,\leq)\text{ be a linearly ordered set and }D\subseteq S.\\
&D\text{ is not order dense in }S
\Longleftrightarrow
\exists x,y\in S\;\bigl(x<y\land \forall d\in D\;(d\leq x\lor y\leq d)\bigr).
\end{aligned}
\]
\[
S\in\mathbf{Set},\qquad D\subseteq S,\qquad P=\mathsf{TotallyOrderedSet}(S,\leq),\qquad
\neg\operatorname{OrderDenseSubset}(D,P)\Longleftrightarrow
\exists x,y\in S\;\bigl(x<y\land \forall d\in D\;(d\leq x\lor y\leq d)\bigr).
\]
\end{description}
\end{remark*}

\begin{remark*}[Interpretation]
Order density is an interval-penetration property. It ignores metric distance
and asks only whether every ordered gap contains a point of the subset. A
subset can therefore be order dense in $S$ precisely when no nonempty open
interval of $S$ is missed entirely by that subset.
\end{remark*}

\begin{dependencies}
\begin{itemize}
  \item \hyperref[def:subset]{Subset}
  \item \hyperref[def:totally-ordered-set]{Total Order}
\end{itemize}
\end{dependencies}
